% Осень26!B8; нумерация внутри столбца без пустых строк.
% Источник: лекция 7, §10 (строка 335)
\definitionticket{06}{Формулировка теоремы о локальной обратимости в терминах систем уравнений}
\label{definition:06}

\textbf{Формулировка той же теоремы в терминах систем уравнений .} Рассмотрим систему \(m\) уравнений с \(m\) неизвестными \[\left\{
\begin{aligned}
f_1(x_1,\dots,x_m)&=y_1,\\
&\ \vdots\\
f_m(x_1,\dots,x_m)&=y_m,
\end{aligned}\right.\] где все \(f_i\) --- класса \(C^1\) (уравнения «приличные»). Пусть при правой части \(b=(b_1,\dots,b_m)\) система имеет решение \(a=(a_1,\dots,a_m)\), и в точке \(a\) отличен от нуля \textbf{определитель Якоби} (матрица составлена из градиентов функций как из строк; \(i\) нумерует уравнения, \(j\) --- переменные): \[\det\Bigl(\frac{\partial f_i}{\partial x_j}(a)\Bigr)_{i,j=1}^{m}\ \ne\ 0.\] Тогда \textbf{для всякого \(y\), близкого к \(b\), система имеет решение \(x\), близкое к \(a\)}, и зависимость \(x=x(y)\) --- класса \(C^1\) (каждая координата \(x_i(y_1,\dots,y_m)\) имеет непрерывные частные производные).

\textbf{Связь с теоремой.} Это в точности локальная обратимость для \(T(x):=(f_1(x),\dots,f_m(x))\), записанная покоординатно: разрешимость при близких правых частях --- биективность на окрестности, класс \(C^1\) зависимости \(x\) от \(y\) --- гладкость обратного отображения. \textbf{Аналогия с линейными системами (лектор).} Для линейной системы \(\alpha_{ij}\): если матрица обратима, то решение существует \textbf{при любых} правых частях и вычисляется по формулам Крамера --- «просто определители, многочлены, знаменатель в ноль не обращается», зависимость \(x(y)\) гладкая (глобально). В нелинейном случае уравнения «вовсе не линейные, но до некоторой степени эффект сохраняется» --- локально.
